% parameter_tables.tex
%
% Landscape A4. Compile with:
%   pdflatex parameter_tables.tex
%   bibtex   parameter_tables
%   pdflatex parameter_tables.tex
%   pdflatex parameter_tables.tex

\documentclass[11pt,a4paper]{article}

\usepackage[landscape, top=2cm, bottom=2cm, left=2cm, right=2cm]{geometry}
\usepackage{booktabs}
\usepackage{longtable}
\usepackage{array}
\usepackage{natbib}
\usepackage{hyperref}
\usepackage{microtype}

% column types
\newcolumntype{L}[1]{>{\raggedright\arraybackslash}p{#1}}
\newcolumntype{C}[1]{>{\centering\arraybackslash}p{#1}}

% suppress page numbers (tables only)
\pagestyle{plain}

% ---- helper macro for the "how we know" column ----
% usage: \conf{text}
\newcommand{\conf}[1]{\textit{#1}}

\begin{document}

\begin{center}
  {\large\bfseries Model Parameter Tables}\\[4pt]
\end{center}

% =============================================================
% TABLE 1 — Know well: fix these
% =============================================================
% Column widths sum to ~25 cm (fits landscape A4 with 4 cm total margin)
% Sym(2) | Desc(5.5) | Value(3) | Units(3.5) | Ref(5) | Basis(6) = 25 cm

\begin{longtable}{L{2cm} L{5.5cm} C{3cm} L{3.5cm} L{5cm} L{6cm}}

\caption{Group~1: well constrained parameters, fixed in all runs.}
\label{tab:known} \\

\toprule
\textbf{Symbol} &
\textbf{Description} &
\textbf{Value} &
\textbf{Units} &
\textbf{Reference} &
\textbf{Notes} \\
\midrule
\endfirsthead

\multicolumn{6}{l}{\small\textit{Table~\ref{tab:known} continued\ldots}} \\[2pt]
\toprule
\textbf{Symbol} &
\textbf{Description} &
\textbf{Value} &
\textbf{Units} &
\textbf{Reference} &
\textbf{Notes} \\
\midrule
\endhead

\midrule
\multicolumn{6}{r}{\small\textit{continued on next page}} \\
\endfoot

\bottomrule
\endlastfoot

% --- Physical constants ---
$\rho_f$ &
Seawater (fluid) density &
1.0275 &
g\,cm$^{-3}$ &
\citet{millero1981} &
\conf{Standard seawater. Millero \& Poisson provide the international equation of state; value changes $<0.1\%$ over North Atlantic conditions.} \\[4pt]

$\nu$ &
Kinematic viscosity &
0.01 &
cm$^{2}$\,s$^{-1}$ &
\citet{sharqawy2010} &
\conf{Seawater at 20\,\textdegree C. Varies $\sim$30\% between 5 and 25\,\textdegree C; currently held fixed (temperature scaling for future work).} \\[4pt]

$g$ &
Gravitational acceleration &
980 &
cm\,s$^{-2}$ &
--- &
\conf{Standard value; invariant for ocean column depths of interest.} \\[4pt]

$k_B$ &
Boltzmann constant &
$1.3\times10^{-16}$ &
erg\,K$^{-1}$ &
--- &
\conf{NIST fundamental constant; enters only the Brownian kernel.} \\[4pt]

% --- Settling law ---
$a_w$ &
Settling speed coefficient (kriest\_8) &
66 &
m\,day$^{-1}$\,cm$^{-b_w}$ &
\citet{kriest2002} &
\conf{Power-law fit to in-situ Atlantic settling data (Eq.~8 of Kriest 2002). Coefficient and exponent chosen together; treat as a pair.} \\[4pt]

$b_w$ &
Settling speed exponent (kriest\_8) &
0.62 &
--- &
\citet{kriest2002} &
\conf{As above. Exponent is robust across datasets; coefficient varies by $\sim$factor 2 between cruises.} \\[4pt]

% --- Fecal pellet sinking ---
$\Delta\rho_{fp}$ &
Fecal pellet excess density &
0.15 &
g\,cm$^{-3}$ &
\citet{turner2002}; \citet{ploug2008} &
\conf{Dense compact pellets ($\sim$150\,kg\,m$^{-3}$ above seawater). Gives $\sim$69\,m\,day$^{-1}$ at bin~8 via Stokes law, consistent with Turner (2002) review and Ploug et al.\ (2008) measurements.} \\[4pt]

% --- Size grid ---
$d_0$ &
Monomer (unit particle) diameter &
$20\times10^{-4}$ &
cm &
\citet{burd2009} &
\conf{20\,\textmu m monomer, typical phytoplankton cell. Sets the size grid; changes shift all bin boundaries.} \\[4pt]

% --- D_max ----
$A_D$ &
$D_\mathrm{max}$ amplitude in $D_\mathrm{max} = A_D \,\varepsilon^{-1/4}$ &
$9.39\times10^{-6}$ &
m$^{3/4}$\,W$^{-1/4}$ &
\citet{burd2009} &
\conf{From Kolmogorov turbulence theory. $D_\mathrm{max}$ sets the largest stable aggregate size per depth layer. Amplitude is better constrained than the exponent.} \\[4pt]

\end{longtable}

\clearpage

% =============================================================
% TABLE 2 — Kind of okay: fit second
% =============================================================

\begin{longtable}{L{2cm} L{5.5cm} C{3cm} L{3.5cm} L{5cm} L{6cm}}

\caption{Group~2: parameters with reasonable literature estimates but site-dependent
         values. Held at literature defaults in the first optimization pass; fit in
         a second pass.}
\label{tab:medium} \\

\toprule
\textbf{Symbol} &
\textbf{Description} &
\textbf{Value} &
\textbf{Units} &
\textbf{Reference} &
\textbf{Basis / confidence} \\
\midrule
\endfirsthead

\multicolumn{6}{l}{\small\textit{Table~\ref{tab:medium} continued\ldots}} \\[2pt]
\toprule
\textbf{Symbol} &
\textbf{Description} &
\textbf{Value} &
\textbf{Units} &
\textbf{Reference} &
\textbf{Basis / confidence} \\
\midrule
\endhead

\midrule
\multicolumn{6}{r}{\small\textit{continued on next page}} \\
\endfoot

\bottomrule
\endlastfoot

% --- Fractal dimension ---
$D_f$ &
Particle fractal dimension &
2.33 &
--- &
\citet{logan1990}; \citet{li1995} &
\conf{Ocean aggregates: range 1.7--2.5 across studies. Value 2.33 is a mid-range choice consistent with marine snow measurements. Controls mass--size relationship and kernel collision cross-sections.} \\[4pt]

% --- Interaction radius ---
$r/r_g$ &
Ratio of collision radius to radius of gyration &
1.6 &
--- &
\citet{li1997} &
\conf{Sets effective collision cross-section for shear and differential settling kernels. Li \& Logan (1997) derive this from fractal geometry; value varies $\sim$1.3--2.0.} \\[4pt]

% --- Zoo: filter feeder ---
$Z_c$ &
Filter feeder (macrozooplankton) concentration &
0.307 &
ind\,m$^{-3}$ &
\citet{stemmann2004a}; \citet{stemmann2004b} &
\conf{Atlantic maximum from Stemmann et al.\ (2004) Table~2. Depth profile from Fig.~1 wired into \texttt{DepthProfile.typical()}. Needs EXPORTS net tow data (Amy / Debbie) for refinement.} \\[4pt]

$c$ &
Filter feeder clearance rate &
0.025 &
m$^{3}$\,ind$^{-1}$\,day$^{-1}$ &
\citet{stemmann2004a} &
\conf{From Stemmann et al.\ (2004). Order-of-magnitude known; exact value uncertain.} \\[4pt]

% --- Zoo: flux feeder ---
$Z_f$ &
Flux feeder concentration &
0.063 &
ind\,m$^{-3}$ &
\citet{stemmann2004a}; \citet{stemmann2004b} &
\conf{Atlantic maximum from Stemmann et al.\ (2004) Table~2. Same depth profile logic as $Z_c$.} \\[4pt]

$s$ &
Flux feeder capture cross-section &
$1.3\times10^{-5}$ &
m$^{2}$\,ind$^{-1}$ &
\citet{stemmann2004a} &
\conf{Stemmann et al.\ (2004). Interpreted as the projected area presented by a sinking particle to a flux feeder.} \\[4pt]

% --- Egestion fraction ---
$p$ &
Egestion fraction (fecal / ingested) &
0.3--0.5 &
--- &
\citet{stemmann2004a} &
\conf{Range reflects food quality and organism type. Default 0.5 (current). Controls fecal pellet production; linked to carbon budget closure.} \\[4pt]

% --- Microbial ---
$r_0$ &
Microbial remineralization base rate &
0.03 &
day$^{-1}$ &
\citet{iversen2013} &
\conf{Iversen \& Ploug (2013) measured respiration on sinking diatom aggregates in a Norwegian fjord. Value represents warm surface conditions; deep column rates $\sim$10$\times$ lower after $Q_{10}$ scaling.} \\[4pt]

$Q_{10}$ &
Temperature sensitivity of microbial rate &
2.0 &
--- &
\citet{iversen2013} &
\conf{Standard marine bacteria value. Doubles the rate per 10\,\textdegree C increase. Widely used; uncertainty $\pm0.5$.} \\[4pt]

% --- Mining ---
$s_m$ &
Micro-zooplankton mining capture cross-section &
$1.3\times10^{-5}$ &
m$^{2}$\,ind$^{-1}$ &
\citet{stemmann2004a} &
\conf{From Stemmann et al.\ (2004) Part I Eq.~25; same value as flux feeder cross-section.} \\[4pt]

% --- Disagg redistribution ---
$f_\mathrm{next}$ &
Fraction of disaggregated biovolume sent to next smaller bin &
2/3 &
--- &
\citet{burd2009} &
\conf{Sectional model redistribution convention. Remainder spread uniformly across smaller bins. No direct measurement; value is a modeling assumption.} \\[4pt]

\end{longtable}

\clearpage

% =============================================================
% TABLE 3 — No idea: fit first
% =============================================================

\begin{longtable}{L{2cm} L{5.5cm} C{3cm} L{3.5cm} L{5cm} L{6cm}}

\caption{Group~3: parameters with no reliable direct field estimate.
         These are fit first against EXPORTS UVP particle size distributions
         (North Atlantic station, surface to 500\,m).
         Grid search over each parameter individually; conjugate gradient or
         simulated annealing for joint optimization if needed.}
\label{tab:unknown} \\

\toprule
\textbf{Symbol} &
\textbf{Description} &
\textbf{Current value} &
\textbf{Units} &
\textbf{Reference} &
\textbf{Why unknown / fit range} \\
\midrule
\endfirsthead

\multicolumn{6}{l}{\small\textit{Table~\ref{tab:unknown} continued\ldots}} \\[2pt]
\toprule
\textbf{Symbol} &
\textbf{Description} &
\textbf{Current value} &
\textbf{Units} &
\textbf{Reference} &
\textbf{Why unknown / fit range} \\
\midrule
\endhead

\midrule
\multicolumn{6}{r}{\small\textit{continued on next page}} \\
\endfoot

\bottomrule
\endlastfoot

% --- Stickiness ---
$\alpha$ &
Collision efficiency (stickiness) of marine snow aggregates &
1.0 &
--- &
\citet{jackson1990}; \citet{burd2009} &
\conf{No direct measurement for natural marine particles. Jackson (1990) set $\alpha=1$ (upper bound). TEP-coated aggregates may have $\alpha \sim 0.01$--$1.0$; value depends on surface chemistry and particle type. \textbf{Fit first}: grid search $\alpha \in [0.1, 1.0]$, step 0.1.} \\[6pt]

% --- Cross-coagulation stickiness ---
$\alpha_\mathrm{cross}$ &
Stickiness between fecal pellets and marine snow &
0.5 &
--- &
--- &
\conf{No field or lab measurement exists for fecal pellet -- marine snow collisions. Fecal pellets are compact ($D_f \approx 2.8$) and low in TEP, so assumed less sticky than aggregates. Current 0.5 is a conservative starting value. \textbf{Fit after} $\alpha$; range [0.1, 1.0].} \\[6pt]

% --- Mining bite size ---
$\delta m$ &
Biovolume removed per encounter (micro-zooplankton mining) &
$10^{-5}$ &
cm$^{3}$ &
\citet{stemmann2004a} &
\conf{Proxy for gut volume of small copepods (\textit{Oncaea} spp.). Stemmann et al.\ (2004) Part I Eq.~25 gives the functional form but does not constrain $\delta m$ directly. Range $10^{-6}$ to $10^{-4}$\,cm$^{3}$. Needs EXPORTS gut-content data (Amy / Debbie). \textbf{Fit with} $\alpha_\mathrm{cross}$.} \\[6pt]

% --- Miner concentration ---
$Z_m$ &
Micro-zooplankton miner concentration &
250 &
ind\,m$^{-3}$ &
\citet{stemmann2004a} &
\conf{Approximate value from Stemmann et al.\ (2004) Fig.~1. Depth profile $Z_m(z)$ is currently set uniform; needs EXPORTS net tow data to set the actual profile. Uncertainty large ($\times 2$--$5$).} \\[6pt]

\end{longtable}

\vspace{10pt}
\noindent\textbf{Optimization note.}
The loss function is
$\mathcal{L} = \sum_{z,\,k} \bigl[\log_{10} N_\mathrm{model}(z,k) - \log_{10} N_\mathrm{obs}(z,k)\bigr]^{2}$,
where $k$ indexes size bin and $z$ indexes depth layer.
Log-space weighting ensures that small and large particles contribute equally.
Observations: EXPORTS processed UVP particle size distributions, North Atlantic station.
See \texttt{scripts/optimize/run\_alpha\_grid\_search.m} for the grid-search implementation
and \texttt{scripts/optimize/loss\_size\_dist.m} for the loss function.

\vspace{10pt}

\bibliographystyle{plainnat}
\bibliography{references}

\end{document}
